\documentclass[aspectratio=169,11pt]{beamer}

\usepackage[utf8]{inputenc}
\usepackage[T1]{fontenc}
\usepackage[brazil]{babel}
\usepackage{lmodern}
\usepackage{xcolor}
\usepackage{booktabs}
\usepackage{tikz}
\usepackage{listings}

\usetikzlibrary{arrows.meta,positioning}

\definecolor{pataGreen}{HTML}{0F766E}
\definecolor{pataDark}{HTML}{102A2A}
\definecolor{pataBlue}{HTML}{2563EB}
\definecolor{pataGray}{HTML}{F3F4F6}
\definecolor{codeBg}{HTML}{F8FAFC}
\definecolor{codeBorder}{HTML}{CBD5E1}

\usetheme{Madrid}
\usecolortheme{default}
\setbeamercolor{palette primary}{bg=pataGreen,fg=white}
\setbeamercolor{palette secondary}{bg=pataDark,fg=white}
\setbeamercolor{palette tertiary}{bg=pataBlue,fg=white}
\setbeamercolor{frametitle}{bg=pataGreen,fg=white}
\setbeamercolor{title}{fg=pataDark}
\setbeamercolor{block title}{bg=pataGreen,fg=white}
\setbeamercolor{block body}{bg=pataGray,fg=black}
\setbeamertemplate{navigation symbols}{}
\setbeamertemplate{itemize item}{\color{pataGreen}$\blacktriangleright$}

\lstdefinestyle{codigo}{
  backgroundcolor=\color{codeBg},
  basicstyle=\ttfamily\scriptsize,
  breaklines=true,
  frame=single,
  rulecolor=\color{codeBorder},
  columns=fullflexible,
  keepspaces=true,
  showstringspaces=false,
  tabsize=2
}

\title[Sistema Pata Feliz]{Sistema de Informação Pata Feliz}
\subtitle{Sistema Web para Clínica Veterinária com PostgreSQL}
\author[Gabriel, Joana, Matheus, Thais e Vicente]{
  Gabriel de Jesus Moraes Mauricio\\
  Joana de Oliveira Rocha\\
  Matheus de Moura Almeida Neri\\
  Thais Pereira Carvalho\\
  Vicente Alves Cordeiro Junior
}
\institute{
  Banco de Dados 2\\
  Professor: Áthila Rocha Trindade
}
\date{2026}

\begin{document}

\begin{frame}
  \titlepage
\end{frame}

\begin{frame}{Objetivo do Sistema}
  \begin{block}{Proposta}
    Desenvolver um Sistema de Informação web para gestão de uma clínica veterinária, com persistência em banco de dados relacional e operações SQL integradas ao backend e ao frontend.
  \end{block}

  \begin{itemize}
    \item Organizar tutores, animais, veterinários, procedimentos e atendimentos.
    \item Controlar agenda, bloqueios, prontuários e lembretes automáticos.
    \item Demonstrar conceitos de Banco de Dados 2 na aplicação real.
  \end{itemize}
\end{frame}

\begin{frame}{Perfis de Usuário}
  \begin{columns}[T]
    \begin{column}{0.32\textwidth}
      \begin{block}{Administrador}
        Gerencia clientes, animais, veterinários, procedimentos, agenda, bloqueios, lembretes e relatórios.
      \end{block}
    \end{column}
    \begin{column}{0.32\textwidth}
      \begin{block}{Veterinário}
        Visualiza agenda, acessa dados do animal e tutor, altera status, registra prontuários e valores adicionais.
      \end{block}
    \end{column}
    \begin{column}{0.32\textwidth}
      \begin{block}{Tutor}
        Cadastra animais, agenda atendimentos, acompanha consultas, visualiza prontuários permitidos e lembretes.
      \end{block}
    \end{column}
  \end{columns}
\end{frame}

\begin{frame}{Requisitos Funcionais}
  \begin{itemize}
    \item Login por e-mail ou CPF e senha.
    \item Controle de permissões por perfil.
    \item Cadastro e edição de clientes, animais, veterinários e procedimentos.
    \item Agendamento de atendimentos com validação de data, horário e disponibilidade.
    \item Registro de prontuários vinculados ao atendimento, animal e veterinário.
    \item Geração de lembretes automáticos para consultas, vacinas e animais sem atendimento recente.
    \item Relatórios de serviços, veterinários e histórico do animal.
  \end{itemize}
\end{frame}

\begin{frame}{Arquitetura}
  \centering
  \begin{tikzpicture}[
    box/.style={draw, rounded corners, fill=pataGray, minimum width=3.5cm, minimum height=1.15cm, align=center},
    arrow/.style={-{Latex[length=3mm]}, thick}
  ]
    \node[box] (front) {Frontend\\HTML, CSS, JavaScript};
    \node[box, right=1.1cm of front] (api) {Backend\\Node.js + Express};
    \node[box, right=1.1cm of api] (db) {PostgreSQL\\Tabelas, Views, Functions, Triggers};
    \draw[arrow] (front) -- node[above]{HTTP/JSON} (api);
    \draw[arrow] (api) -- node[above]{SQL} (db);
    \draw[arrow] (db) -- node[below]{Registros} (api);
    \draw[arrow] (api) -- node[below]{JSON} (front);
  \end{tikzpicture}

  \vspace{0.7cm}
  \begin{itemize}
    \item Frontend consome a API.
    \item Backend aplica autenticação, autorização e regras de negócio.
    \item Banco mantém integridade, consultas, funções e automações.
  \end{itemize}
\end{frame}

\begin{frame}{Modelo de Dados}
  \begin{columns}[T]
    \begin{column}{0.48\textwidth}
      \begin{itemize}
        \item \texttt{usuarios}
        \item \texttt{clientes}
        \item \texttt{veterinarios}
        \item \texttt{animais}
        \item \texttt{procedimentos}
        \item \texttt{atendimentos}
      \end{itemize}
    \end{column}
    \begin{column}{0.48\textwidth}
      \begin{itemize}
        \item \texttt{prontuarios}
        \item \texttt{vacinas\_aplicadas}
        \item \texttt{lembretes}
        \item \texttt{bloqueios\_agenda}
        \item \texttt{historico\_precos}
        \item \texttt{historico\_pesos}
      \end{itemize}
    \end{column}
  \end{columns}

  \begin{block}{Requisito atendido}
    O sistema manipula dados armazenados em mais de uma estrutura de armazenamento, com múltiplas tabelas relacionadas por chaves estrangeiras.
  \end{block}
\end{frame}

\begin{frame}{CRUD no Sistema}
  \begin{table}
    \centering
    \begin{tabular}{lll}
      \toprule
      Comando & Exemplo no sistema & Finalidade \\
      \midrule
      \texttt{SELECT} & Listagens e relatórios & Consultar dados \\
      \texttt{INSERT} & Cadastro de tutores e atendimentos & Inserir dados \\
      \texttt{UPDATE} & Edição de perfil e status & Atualizar dados \\
      \texttt{DELETE} & Exclusão de lembretes e bloqueios & Remover dados \\
      \bottomrule
    \end{tabular}
  \end{table}

  \begin{block}{Integração}
    As operações são chamadas pelo backend Express e persistidas no PostgreSQL.
  \end{block}
\end{frame}

\begin{frame}[fragile]{Trecho de Código: Conexão com o Banco}
  \begin{lstlisting}[style=codigo,language=JavaScript]
const pool = new Pool({
  host: process.env.DB_HOST,
  user: process.env.DB_USER,
  password: process.env.DB_PASSWORD,
  database: process.env.DB_NAME,
  port: process.env.DB_PORT || 5432
});

await pool.query(sql, params);
  \end{lstlisting}

  \begin{itemize}
    \item Trecho localizado em \texttt{backend/db.js}.
    \item Uso de variáveis de ambiente para credenciais.
    \item Biblioteca \texttt{pg} para PostgreSQL.
  \end{itemize}
\end{frame}

\begin{frame}[fragile]{Transações}
  \begin{lstlisting}[style=codigo,language=JavaScript]
await conn.beginTransaction();

await conn.query("INSERT INTO usuarios (...) VALUES (...)");
await conn.query("INSERT INTO clientes (...) VALUES (...)");

await conn.commit();
  \end{lstlisting}

  \begin{block}{Onde foi usado}
    Cadastro e edição de clientes, veterinários e perfil, pois envolvem mais de uma tabela.
  \end{block}
\end{frame}

\begin{frame}[fragile]{Trecho de Código: Consulta com Junções}
  \begin{lstlisting}[style=codigo,language=SQL]
SELECT
  at.id_atendimento,
  an.nome AS animal,
  uc.nome AS tutor,
  uv.nome AS veterinario,
  p.nome AS procedimento
FROM atendimentos at
JOIN animais an ON an.id_animal = at.id_animal
JOIN clientes c ON c.id_cliente = at.id_cliente
JOIN usuarios uc ON uc.id_usuario = c.id_usuario
JOIN veterinarios v ON v.id_veterinario = at.id_veterinario
JOIN usuarios uv ON uv.id_usuario = v.id_usuario
JOIN procedimentos p ON p.id_procedimento = at.id_procedimento;
\end{lstlisting}
\begin{itemize}
  \item Consulta usada na agenda e nas funções de listagem.
  \item Relaciona atendimentos, animais, tutores, veterinários e procedimentos.
\end{itemize}
\end{frame}

\begin{frame}[fragile]{Trecho de Código: Procedure e Functions}
  \begin{itemize}
    \item Procedure real: \texttt{atualizar\_status\_lembrete(id, status, usuario)}
    \item \texttt{buscar\_historico\_animal(id)}
    \item \texttt{listar\_clientes()}
    \item \texttt{listar\_animais(tipo, cliente)}
    \item \texttt{listar\_atendimentos(tipo, cliente, veterinario)}
    \item \texttt{listar\_prontuarios(tipo, cliente, veterinario)}
    \item \texttt{listar\_lembretes(tipo, cliente, veterinario)}
    \item \texttt{relatorio\_servicos()} e \texttt{relatorio\_veterinarios()}
  \end{itemize}

  \begin{lstlisting}[style=codigo,language=JavaScript]
db.query(
  "CALL atualizar_status_lembrete(?::bigint, ?::status_lembrete, ?::bigint)",
  [idLembrete, status, idUsuario]
);

db.query(
  "SELECT * FROM listar_atendimentos(?::tipo_usuario, ?::bigint, ?::bigint)",
  [tipoUsuario, idCliente, idVeterinario]
);
  \end{lstlisting}
\begin{itemize}
  \item Chamada da procedure em \texttt{backend/routes/v2.js}.
  \item Functions usadas para listagens e relatórios do sistema.
\end{itemize}
\end{frame}

\begin{frame}[fragile]{Triggers e Automações}
  \begin{lstlisting}[style=codigo,language=SQL]
CREATE TRIGGER trg_atendimentos_gerar_lembretes
AFTER INSERT OR UPDATE OF inicio, fim, status,
  id_animal, id_cliente, id_procedimento
ON atendimentos
FOR EACH STATEMENT
EXECUTE FUNCTION disparar_geracao_lembretes();
  \end{lstlisting}

  \begin{itemize}
    \item Atualização automática do campo \texttt{atualizado\_em}.
    \item Geração automática de lembretes após mudanças em atendimentos e vacinas.
  \end{itemize}
\end{frame}

\begin{frame}{Regras de Agendamento}
  \begin{itemize}
    \item Não permite atendimento no passado.
    \item Não permite agendamento com mais de um ano de antecedência.
    \item Segunda a sexta: 08:00 às 18:00.
    \item Sábado: 08:00 às 12:00.
    \item Domingo indisponível.
    \item Bloqueios de agenda impedem novos atendimentos no período.
    \item Restrição \texttt{EXCLUDE USING GIST} evita choque para o mesmo veterinário.
  \end{itemize}
\end{frame}

\begin{frame}{Demonstração}
  \begin{enumerate}
    \item Login no sistema com usuário de teste.
    \item Visualização do dashboard por perfil.
    \item Cadastro ou edição de animal/procedimento.
    \item Criação de atendimento na agenda.
    \item Registro de prontuário.
    \item Visualização de lembretes e relatórios.
  \end{enumerate}

  \begin{block}{Dados de exemplo}
    Usuários fictícios foram inseridos via seed PostgreSQL. Senha padrão: \texttt{123456}.
  \end{block}
\end{frame}

\begin{frame}{Como Executar}
  \begin{itemize}
    \item Criar banco PostgreSQL: \texttt{pata\_feliz\_v2}.
    \item Executar \texttt{schema-v2-postgres.sql}.
    \item Executar \texttt{seed-v2-postgres.sql}.
    \item Atualizar funções com \texttt{functions-v2-postgres.sql}, se necessário.
    \item Rodar backend com \texttt{npm start}.
    \item Abrir o frontend no navegador.
  \end{itemize}
\end{frame}

\begin{frame}{Conclusão}
  \begin{block}{Requisitos atendidos}
    O Pata Feliz implementa um Sistema de Informação web funcional com banco relacional, CRUD completo, transações, consultas com junções, funções armazenadas com parâmetros, triggers e múltiplas tabelas relacionadas.
  \end{block}

  \begin{itemize}
    \item O sistema resolve uma necessidade prática de clínica veterinária.
    \item O banco de dados participa das regras de integridade e automação.
    \item A aplicação demonstra integração entre frontend, backend e PostgreSQL.
  \end{itemize}
\end{frame}

\begin{frame}
  \centering
  {\Huge Obrigado!}

  \vspace{0.8cm}
  {\Large Dúvidas?}
\end{frame}

\end{document}
